%============================================================
% CH1 · PATTERN 2 — EXERCISES, HINTS, ANSWERS
%============================================================
\subsection{Pattern 2: Exercises and Going Further}
\label{ssec:ch1-p2-ex}

\begin{patternbox}{Pattern 2 Summary Card}
\begin{tabularx}{\linewidth}{@{}>{\bfseries}l X@{}}
Recognise & ``remove / move / compact in place'', ``return the new length'', ``keep the relative order'', $\bigO{1}$ extra memory \\
Core loop & \texttt{w = 0; for r in [0, n): if keep(a[r], a[0..w)) a[w++] = a[r];} \\
Invariant & $w \le r$; $a[0..w)$ = kept elements of $a[0..r)$ in order; $a[r..n)$ untouched \\
Keep rules & the element alone; the last kept elements ($a_{w-1}$, $a_{w-k}$); cancellation ($w{-}{-}$); never the cells in $[w, r)$ \\
Mirror & output longer than input: both pointers from the right, $w \ge r$ \\
Top traps & \texttt{erase} in a loop; comparing with $a_{r-k}$; forgetting the fill or the returned length; a forward pointer for insertions \\
\end{tabularx}
\end{patternbox}

\exlegend
\begin{enumerate}[leftmargin=*, label=\textbf{E2.\arabic*}, itemsep=6pt]
    \item \exeasy\ \ptricky \exprob{Duplicate Zeros}{LeetCode 1089}{https://leetcode.com/problems/duplicate-zeros/}\ \lc{Easy}. Duplicate every zero of a fixed-length array in place, shifting the other elements to the right; elements pushed beyond the end are lost. $n \le 10^4$, $0 \le a_i \le 9$.
    \item \exeasy\ \pclassic \exprob{Restore the Permutation by Merger}{Codeforces 1385B}{https://codeforces.com/problemset/problem/1385/B}\ \cf{800}. A permutation $p$ of $1..n$ was merged with a copy of itself (keeping the relative order within each copy), giving a sequence of length $2n$. Restore $p$. $n \le 50$, $t \le 400$.
    \item \exeasy\ \precent \exprob{Clear Digits}{LeetCode 3174}{https://leetcode.com/problems/clear-digits/}\ \lc{Easy}. A string of lowercase letters and digits is given. Repeatedly delete the first digit together with the closest non-digit character to its left. Return the final string. $|s| \le 100$; the input guarantees every digit can be deleted.
    \item \exmedium\ \pconcept \exprob{Rotating the Box}{LeetCode 1861}{https://leetcode.com/problems/rotating-the-box/}\ \lc{Medium}. An $m \times n$ grid holds stones \texttt{\#}, fixed obstacles \texttt{*} and empty cells \texttt{.}. Rotate the box $90^\circ$ clockwise; gravity then makes every stone fall until it lands on an obstacle, another stone or the bottom. Return the $n \times m$ result. $m, n \le 500$.
    \item \exmedium\ \ptricky \exown{Normalise separators.} An array of integers uses $0$ as a separator between blocks of non-zero values, but the zeros are messy: several zeros in a row, zeros at the start and at the end. Rewrite it in place so that consecutive blocks are separated by exactly one zero and there are no leading or trailing zeros; return the new length. Example: $[0,0,3,4,0,0,0,5,0,6,0,0] \to [3,4,0,5,0,6]$, length 6.
    \item \exmedium\ \pconcept \exown{Delete a list of positions.} You are given an array $a$ of length $n$ and a \emph{sorted} list of $m$ distinct positions $p_1 < \dots < p_m$. Delete all these positions from $a$ in place in $\bigO{n + m}$ time and $\bigO{1}$ extra memory, keeping the order of the rest; return the new length.
\end{enumerate}

\subsubsection*{Hints}
\begin{enumerate}[leftmargin=*, label=\textbf{E2.\arabic*}]
    \item A forward write pointer would run ahead of the read pointer. Where does the last element that survives come from? Fill the array from the right.
    \item The first time a value appears, its first copy has just been read. A boolean array \texttt{seen[1..n]} decides which elements to keep.
    \item A digit removes the last kept letter: the write pointer moves \emph{back}.
    \item Rotation and gravity can be separated: in the original grid, gravity after the rotation pulls stones to the \emph{right} within each row. An obstacle resets where the next stone can land.
    \item Decide each zero by looking at the last \emph{written} element. What can the last written element be when the loop ends?
    \item Keep a second read pointer $j$ into the list of positions. Is index $r$ the next position to delete?
\end{enumerate}

\subsubsection*{Answers}

\begin{solutionbox}[E2.1 Duplicate Zeros]
The output is longer than the kept input, so the write pointer must not start at the left. Count the zeros $z$: without truncation, element $i$ would move to position $i + (\text{number of zeros before } i)$. Walk $i$ from the right with a virtual write position $j$ that starts at $n - 1 + z$; write only when $j < n$ (cells beyond the end are discarded), and write a zero twice. Since $j \ge i$ always, no unread cell is overwritten --- the mirror image of the invariant $w \le r$.
\begin{lstlisting}[style=cpp]
#include <bits/stdc++.h>
using namespace std;

class Solution {
public:
    void duplicateZeros(vector<int>& arr) {
        int n = arr.size(), zeros = 0;
        for (int x : arr) if (x == 0) zeros++;
        for (int i = n - 1, j = n - 1 + zeros; i >= 0; i--) {   // j: virtual write position
            if (j < n) arr[j] = arr[i];
            if (arr[i] == 0) {                                  // a zero is written twice
                j--;
                if (j < n) arr[j] = 0;
            }
            j--;
        }
    }
};

int main() {
    vector<vector<int>> tests = {{1, 0, 2, 3, 0, 4, 5, 0}, {1, 2, 3}};
    for (auto& t : tests) {
        Solution().duplicateZeros(t);
        for (int x : t) cout << x << " ";      // 1 0 0 2 3 0 0 4 | 1 2 3
        cout << "| ";
    }
    cout << "\n";
}
\end{lstlisting}
On Example 1 ($z = 3$, $j$ starts at 10): $i = 7, 6$ write nothing (virtual positions $8, 9, 10$ fall off), $i = 5$ writes the 4 at index 7, $i = 4$ writes zeros at 6 and 5 --- the state of the last diagram of the pattern theory --- and so on, ending with $[1, 0, 0, 2, 3, 0, 0, 4]$. The forward version overwrites the 2 after the first zero and produces $[1, 0, 0, 0, \dots]$. Checked against building the answer in a new vector on 30000 random arrays.
\end{solutionbox}

\begin{solutionbox}[E2.2 Restore the Permutation by Merger]
The first occurrence of each value in the merged sequence comes from the first copy of $p$ \emph{or} from the second; either way the first occurrences appear in the order of $p$, because both copies list the values in that order and each value's first occurrence precedes its second. So $p$ is the subsequence of first occurrences: a read/write filter with a \texttt{seen} array as the keep rule.
\begin{lstlisting}[style=cpp]
#include <bits/stdc++.h>
using namespace std;

int main() {
    int t;
    cin >> t;
    while (t--) {
        int n;
        cin >> n;
        vector<int> a(2 * n);
        for (int& x : a) cin >> x;
        vector<bool> seen(n + 1, false);
        int w = 0;
        for (int r = 0; r < 2 * n; r++)
            if (!seen[a[r]]) { seen[a[r]] = true; a[w++] = a[r]; }   // keep first occurrences
        for (int i = 0; i < w; i++) cout << a[i] << (i + 1 < w ? ' ' : '\n');
    }
}
\end{lstlisting}
Official sample: $1\ 2$; $1\ 3\ 4\ 2$; $1\ 2\ 3\ 4\ 5$; $1\ 2\ 3$; $2\ 3\ 4\ 1$. The \texttt{seen} array is a frequency array over the small values $1..n$ (Pattern~5) used as a filter.
\end{solutionbox}

\begin{solutionbox}[E2.3 Clear Digits]
The ``closest non-digit to the left'' that is still present is exactly the last kept letter, so the kept prefix is a stack: a letter is pushed (\texttt{s[w++] = c}), a digit pops (\texttt{w--}). The guarantee says a letter is always available; the test \texttt{w > 0} makes the code safe anyway.
\begin{lstlisting}[style=cpp]
#include <bits/stdc++.h>
using namespace std;

class Solution {
public:
    string clearDigits(string s) {
        int w = 0;
        for (char c : s) {
            if (isdigit(c)) { if (w > 0) w--; }   // the digit deletes the last kept letter
            else s[w++] = c;                       // keep the letter
        }
        return s.substr(0, w);
    }
};

int main() {
    cout << "[" << Solution().clearDigits("abc") << "] [" << Solution().clearDigits("cb34")
         << "] [" << Solution().clearDigits("ab1c2d") << "]\n";   // [abc] [] [ad]
}
\end{lstlisting}
On \texttt{ab1c2d}: push \texttt{a}, \texttt{b}; the 1 pops \texttt{b}; push \texttt{c}; the 2 pops \texttt{c}; push \texttt{d}: result \texttt{ad}. The repeated ``find the first digit, delete it and its letter'' simulation is $\bigO{n^2}$ and gives the same results (checked on 30000 random valid strings). Note that deleting the \emph{first} digit first does not matter for the stack view: each digit always removes the letter that is closest to it at that moment.
\end{solutionbox}

\begin{solutionbox}[E2.4 Rotating the Box]
After a clockwise rotation, ``down'' is the direction of increasing column index in the original grid. So: (1) in every row, let the stones slide to the right; (2) rotate. Step (1) is a write pointer per row moving right to left: \texttt{w} is the cell where the next stone lands; an obstacle at column $r$ resets it to $r - 1$; a stone at $r$ is lifted (\texttt{.}) and dropped at \texttt{w--}. Since $w \ge r$, the stone is never dropped on an unread cell that still holds a stone.
\begin{lstlisting}[style=cpp]
#include <bits/stdc++.h>
using namespace std;

class Solution {
public:
    vector<vector<char>> rotateTheBox(vector<vector<char>>& box) {
        int m = box.size(), n = box[0].size();
        for (auto& row : box) {
            int w = n - 1;                              // where the next stone lands
            for (int r = n - 1; r >= 0; r--) {
                if (row[r] == '*') w = r - 1;           // obstacle: a new segment starts left of it
                else if (row[r] == '#') { row[r] = '.'; row[w--] = '#'; }   // slide right
            }
        }
        vector<vector<char>> res(n, vector<char>(m));
        for (int i = 0; i < m; i++)
            for (int j = 0; j < n; j++) res[j][m - 1 - i] = box[i][j];   // rotate clockwise
        return res;
    }
};

int main() {
    vector<vector<char>> box = {{'#', '.', '*', '.'}, {'#', '#', '*', '.'}};
    for (auto& row : Solution().rotateTheBox(box)) {
        for (char c : row) cout << c;
        cout << " ";                                    // #. ## ** ..
    }
    cout << "\n";
}
\end{lstlisting}
\begin{center}
\begin{tikzpicture}[x=0.75cm, y=0.75cm, b/.style={gcell, minimum size=6.5mm, font=\small\ttfamily}]
    \foreach \row [count=\i from 0] in {{\#,.,*,.},{\#,\#,*,.}}{
        \foreach \c [count=\j from 0] in \row { \node[b] at (\j,-\i) {\c}; }
    }
    \node[font=\scriptsize, text=darkgray] at (1.5,0.75) {input};
    \draw[flow] (3.6,-0.5) -- node[lbl, above]{slide $\rightarrow$} (5.4,-0.5);
    \foreach \row [count=\i from 0] in {{.,\#,*,.},{\#,\#,*,.}}{
        \foreach \c [count=\j from 0] in \row { \node[b] at (6+\j,-\i) {\c}; }
    }
    \draw[flow] (9.6,-0.5) -- node[lbl, above]{rotate} (11.2,-0.5);
    \foreach \row [count=\i from 0] in {{\#,.},{\#,\#},{*,*},{.,.}}{
        \foreach \c [count=\j from 0] in \row { \node[b] at (12+\j,0.75-\i) {\c}; }
    }
\end{tikzpicture}
\end{center}
Example 2: the stone in row 0 slides one cell right and stops at the obstacle; row 1 is already packed. All three official examples are reproduced, and the result equals a step-by-step gravity simulation on 30000 random grids. Time $\bigO{mn}$, extra memory only for the rotated output.
\end{solutionbox}

\begin{solutionbox}[E2.5 Normalise separators]
A non-zero is always kept. A zero is kept only if the last \emph{written} element is a non-zero (so the output never starts with a zero and never has two in a row). The trap is the end: if the input ends with zeros, the loop writes one separator after the last block and nothing follows it --- undo it after the loop.
\begin{lstlisting}[style=cpp]
#include <bits/stdc++.h>
using namespace std;

int normalise(vector<int>& a) {
    int w = 0;
    for (int r = 0; r < (int)a.size(); r++) {
        int x = a[r];
        if (x != 0) a[w++] = x;                        // block element: always kept
        else if (w > 0 && a[w - 1] != 0) a[w++] = 0;   // one separator after a block
    }
    if (w > 0 && a[w - 1] == 0) w--;                   // no trailing separator
    return w;
}

int main() {
    vector<int> a = {0, 0, 3, 4, 0, 0, 0, 5, 0, 6, 0, 0};
    int k = normalise(a);
    cout << k << ":";
    for (int i = 0; i < k; i++) cout << " " << a[i];   // 6: 3 4 0 5 0 6
    cout << "\n";
    vector<int> z = {0, 0, 0};
    cout << normalise(z) << "\n";                      // 0
}
\end{lstlisting}
\begin{center}
\begin{tikzpicture}
    \dparray{n}{0.6}{in}{0/aout, 0/aout, 3/acell, 4/acell, 0/acell, 0/aout, 0/aout, 5/acell, 0/acell, 6/acell, 0/acell, 0/aout}
    \dparray{o}{-0.3}{out}{3/akeep, 4/akeep, 0/akeep, 5/akeep, 0/akeep, 6/akeep}
\end{tikzpicture}
\end{center}
Without the final \texttt{w--} the example returns length 7 with a trailing 0. Keeping every first zero of a run ``because the previous element was non-zero'' while comparing with the \emph{original} neighbour \texttt{a[r-1]} also works here, but only because $a_{r-1}$ is never overwritten before $a_r$ is read; the $a_{w-1}$ rule states directly what the output needs. Checked against ``split into blocks, join with single zeros'' on 50000 random arrays.
\end{solutionbox}

\begin{solutionbox}[E2.6 Delete a list of positions]
Two read pointers and one write pointer: $r$ scans the array, $j$ scans the sorted positions. If $p_j = r$, index $r$ is deleted and $j$ advances; otherwise $a_r$ is copied to $a_{w++}$. Because the positions are sorted, $p_j$ is always the smallest position not yet passed, and every position is checked exactly when $r$ reaches it.
\begin{lstlisting}[style=cpp]
#include <bits/stdc++.h>
using namespace std;

int deletePositions(vector<int>& a, const vector<int>& pos) {   // pos sorted, distinct
    int w = 0, j = 0;
    for (int r = 0; r < (int)a.size(); r++) {
        if (j < (int)pos.size() && pos[j] == r) { j++; continue; }   // r is deleted
        a[w++] = a[r];
    }
    return w;
}

int main() {
    vector<int> a = {10, 20, 30, 40, 50, 60};
    int k = deletePositions(a, {0, 2, 3});
    for (int i = 0; i < k; i++) cout << a[i] << " ";   // 20 50 60
    cout << "\n";
}
\end{lstlisting}
Time $\bigO{n + m}$, memory $\bigO{1}$. Deleting the positions one by one with \texttt{erase} would cost $\bigO{nm}$ --- and would also shift the later positions, so $p_2, p_3, \dots$ would have to be corrected by the number of earlier deletions. Checked against a filter into a new vector on 30000 random inputs.
\end{solutionbox}
